
\section{主要结果的逐章概述}
\label{sec:1-5}

从整体上看，拉格朗日交弗洛尔理论是一套极其复杂而非平凡的理论；部分原因在于，许多基本构造都必须在链层面完成。这些链层面的构造，又同如今已广为人知的各类伪全纯映射模空间之（虚）基本链构造错综交织；后一构造借助仓西结构 \textbf{[FuOn99II]} 或同类的等价工具来完成。这在一定程度上解释了：除了辛几何以外，这套理论为何还需要现代数学中的许多学科，例如无穷维拓扑、非线性椭圆偏微分方程、形变理论以及高阶代数结构。不仅如此，带 D-膜的开弦理论之物理学还进一步丰富了这套理论。我们力图为这一复杂理论的各个方面建立严格基础，同时把它在辛拓扑和同调镜像对称两方面的应用纳入考虑。

这也解释了为何本书作为一部研究专著，篇幅显得异乎寻常地庞大。我们预计，不仅相关领域的数学家，而且研究弦理论的物理学家，也会依照各自的兴趣至少阅读本书的一部分。因此，我们想逐章简要概述全书内容，帮助读者判断应当到哪里寻找自己感兴趣的内容。

下面按章概述本书内容。在此过程中，我们还将简述前面已经陈述的各个主要定理之证明。需要指出的是，许多章都有各自的引言（第 3.1、4.1、6.1 节以及第 7 章开头），其中给出了更详细的本章概要。

第 2 章主要回顾弗洛尔最初版本的弗洛尔上同调 \textbf{[Flo88IV]}，并说明一旦离开弗洛尔在 \textbf{[Flo88IV]} 中所考虑的拉格朗日子流形 $L$ 之范围，这一定义如何立即遇到障碍；这里所说的范围，是指没有量子效应的情形，或者满足
$\pi_2(M,L)=\{e\}$ 的 $L$。第 2.1--2.2 节定义各种马斯洛夫指标，以及本书全部构造所依托的伪全纯圆盘稳定映射模空间。第 2.3 节说明由弗洛尔本人提出的拉格朗日交弗洛尔上同调构造之基本思想 \textbf{[Flo88IV]}。随后，第 2.4 节首先回顾第二署名作者在 \textbf{[Oh93,95I]} 中研究的单调拉格朗日子流形之弗洛尔上同调，再说明这些工作为何不能直接推广到一般情形，并指出可能克服这一失败的办法。

我们基本结果定理 A 和 F 的证明从第 3 章开始。第 3--5 章进行一系列基本构造，以定义我们的形变弗洛尔上同调，继而证明定理 A 和 F。为此，我们使用各类伪全纯圆盘模空间的（虚）基本链。构造所需的两个要点则推迟到第 7、8 章处理：第 7 章解决横截性问题，第 8 章解决定向问题。

第 3--5 章一贯采用如下写法：先在纯代数语境中发展滤过 $A_\infty$ 代数和滤过 $A_\infty$ 模所需的同调（同伦）代数，再把这套代数工具应用于拉格朗日子流形的几何语境。这样一来，同调（同伦）代数部分便可以脱离弗洛尔上同调的几何理论而单独阅读；这对偏重代数的读者或许有所裨益。第 3、4 章和第 5 章大部分几何构造都必须在（余）链层面进行。给定拉格朗日子流形 $L\subset(M,\omega)$，我们在 $L$ 的一个余链复形上构造其滤过 $A_\infty$ 结构；该复形的上同调群与 $L$ 的上同调群同构。与通常的上同调理论一样，实现 $L$ 的上同调群的余链复形有若干种不同选择。在撰写本书期间乃至此前，我们试验过多种选择。最终，本版主要采用奇异链复形。更准确地说，我们在 $L$ 的光滑奇异链复形的一个可数生成子复形上构造滤过 $A_\infty$ 结构；我们把 $k$-链认作 $\dim L-k$ 次余链，从而将其视为余链复形。（这可以看作一种庞加莱对偶。）

第 3 章给出滤过 $A_\infty$ 代数及其形变理论的定义。我们先在代数语境中引入障碍闭链和定界余链的概念，再在第 3.5 节以 $L$ 上一个可数生成复形为底层复形，为每个拉格朗日子流形 $L$ 配置一个滤过 $A_\infty$ 代数。第 3.6 节给出相应的代数论证，说明如何从滤过 $A_\infty$ 代数出发，定义（形变）弗洛尔上同调以及参数化这些弗洛尔上同调的模空间 $\mathcal M(L)$。单个拉格朗日子流形 $L$ 上的形变弗洛尔上同调，应当视作一种博特--莫尔斯型弗洛尔上同调。接着，第 3.7 节为每一对拉格朗日子流形配置一个滤过 $A_\infty$ 双模。定理 A 和 F 的陈述所用到的、直接定义在上同调群上的滤过 $A_\infty$ 结构，直到第 5 章末尾才构造出来。

第 3 章还讨论另外两个重要问题。其一是滤过 $A_\infty$ 代数的幺元。由于第 7 章所解决的横截性问题，我们必须引入同伦幺元的概念；其定义见第 3.3 节。幺元与弱定界余链模空间 $\mathcal M_{\mathrm{weak}}(L)$ 之间的关系，在第 3.6.3 小节说明。

其二是第 3.8 节讨论的形变：用来自环境辛流形的上同调类来形变滤过 $A_\infty$ 代数。为此，我们使用记作 $\mathfrak p$、$\mathfrak q$ 和 $\mathfrak r$ 的若干算子。为了定义这些算子，我们采用在圆盘内部和边界上都带标记点的伪全纯圆盘模空间。这是在弗洛尔理论中纳入开闭弦模空间的第一步。我们预计，开闭弦理论的思想还会更广泛地渗入弗洛尔理论和辛几何（见 \textbf{[FOOO08II]}）。本书证明定理 H 和 I 时已经运用了这样的思想。算子 $\mathfrak q$ 与模空间 $\mathcal M_{\mathrm{def}}(L)$ 的关系在第 3.8 节讨论，而第 6 章则用算子 $\mathfrak p$ 研究谱序列并证明 (D.3)。

第 3 章在奇异余链复形的可数生成子复形上构造出的滤过 $A_\infty$ 代数，依赖于构造中所作的各种选择。第 4、5 章讨论它在何种意义下与这些选择无关，并证明滤过 $A_\infty$ 代数在滤过 $A_\infty$ 代数的同伦等价意义下唯一。不过，在陈述这一同伦等价之前，我们需要定义 $A_\infty$ 结构同调代数中的若干基本概念。第 4 章定义（滤过）$A_\infty$ 代数之间的同伦等价，并证明其基本性质；本章还详细阐述（滤过）$A_\infty$ 代数的同调（同伦）代数。文献中已经有一些地方定义了 $A_\infty$ 代数之间的同伦等价（例如 \textbf{[Smi00, MSS02]}），但我们仍尽可能给出自足的叙述。这有若干原因。

原因之一是，文献采用了数种不同版本的同伦等价定义；我们认为，如果为了本书的目的而直接借用它们，会造成太多混淆。定义 4.2.35 给出我们对两个 $A_\infty$ 同态之间同伦的定义。我们的想法，是采用一个概念作为几何中“空间与区间 $[0,1]$ 作乘积”这一操作的代数对应物，并在定义 $A_\infty$ 同态的同伦时模仿这一操作。依照在代数上实现“乘以区间 $[0,1]$”的方式不同，会得到数种不同版本的同伦定义；不过，我们证明了这些不同选择在精确意义下全都等价。（这是定理 4.2.34 的一个推论。）

叙述保持自足的另一个原因，是我们还需要沿两个不同方向推广 $A_\infty$ 代数的同调代数：一是定义 $A_\infty$ 结构的滤过版本，二是考虑系数环不是 $\mathbb Q$、而是 $\mathbb Z$ 或 $\mathbb Z_2$ 的情形。我们的叙述经过安排，使这种推广能够容易地完成。

下面介绍滤过 $A_\infty$ 结构同调代数中的两个基本结果。设
\[
  \mathfrak f\colon(C_1,\mathfrak m_*)
  \longrightarrow(C_2,\mathfrak m_*)
\]
是环 $\Lambda_{0,\mathrm{nov}}(R)$ 上两个滤过 $A_\infty$ 代数之间的滤过 $A_\infty$ 同态。假定 $C_i$ 和 $\mathfrak f$ 在定义 3.2.26、3.2.29 的意义下都是有隙的。把系数环 $\Lambda_{0,\mathrm{nov}}(R)$ 约化到 $R$，便得到 $R$ 上的 $A_\infty$ 同态
\[
  \overline{\mathfrak f}\colon
  (\overline C_1,\overline{\mathfrak m}_*)
  \longrightarrow
  (\overline C_2,\overline{\mathfrak m}_*).
\]
对于（非滤过）$A_\infty$ 代数 $(\overline C,\overline{\mathfrak m}_*)$，我们总是假定 $\overline{\mathfrak m}_0=0$，于是
$\overline{\mathfrak m}_1\circ\overline{\mathfrak m}_1=0$。换言之，在非滤过情形中，我们总是假定有一个真正的余链复形 $(\overline C,\overline{\mathfrak m}_1)$。在我们的几何应用中，它无非是 $L$ 的某个奇异链复形版本。

现在陈述滤过 $A_\infty$ 代数同调代数中的第一个主要定理；它是拓扑学中
Whitehead 定理的代数对应物。

\begin{theorem}{V}\label{thm:intro-V}
若 $\overline{\mathfrak f}$ 诱导链同伦等价
\[
  (\overline C_1,\overline{\mathfrak m}_1)
  \longrightarrow
  (\overline C_2,\overline{\mathfrak m}_1),
\]
则
\[
  \mathfrak f\colon(C_1,\mathfrak m_*)
  \longrightarrow(C_2,\mathfrak m_*)
\]
有一个同伦逆（在有隙滤过 $A_\infty$ 代数的意义下）。
\end{theorem}

定理 V 可用来证明，本书各处定义的许多滤过 $A_\infty$ 同态都是同伦等价。需要指出的是，在微分分次代数范畴中，定理 V 一般并不成立。由于定理 V，我们无须“逆转拟同构”；而后者是微分分次代数的标准同调（同伦）代数中的一个关键环节。正因如此，$A_\infty$ 代数的同伦论比微分分次代数的同伦论更接近空间的同伦论。

定理 V 就是定理 4.2.45，其证明见第 4.5 节。接着，第 4.6 节利用这套代数工具，证明第 3 章为拉格朗日子流形 $L$ 构造的滤过 $A_\infty$ 代数在上述同伦等价意义下与各种选择无关。

另一个需要使用 $A_\infty$ 代数同调代数的问题，是式 (1.8) 的解之间的规范等价。我们把式 (1.8) 的解之规范等价类所成的集合定义为 $\mathcal M(L)$。两个解之间的规范等价被定义为一种同伦；它同样是“与 $[0,1]$ 作乘积”的代数对应物。这样一来，就能以函子方式为每个滤过 $A_\infty$ 代数 $(C,\mathfrak m_*)$ 配置集合 $\mathcal M(C)$（定理 4.3.13），继而证明 $\mathcal M(C)$ 是滤过 $A_\infty$ 代数 $(C,\mathfrak m_*)$ 的同伦型不变量（推论 4.3.14）。

第 5 章给出第 4 章的滤过 $A_\infty$ 双模版本。第 5.1 节首先讨论 $\Lambda_{0,\mathrm{nov}}$ 等各种诺维科夫环之间的关系，以及它们在拉格朗日子流形移动时如何变化。第 5.2 节讨论滤过 $A_\infty$ 双模的同调代数，尤其定义它们之间的同伦。随后，第 5.3 节证明，为拉格朗日子流形的相对 spin 对 $(L_0,L_1)$ 配置的滤过 $A_\infty$ 双模，在链同伦等价意义下与构造中涉及的各种选择无关。

第 5.4 节把在余链层面构造的各种 $A_\infty$ 结构转移到上同调上。就此而言，下面是滤过 $A_\infty$ 结构同调代数中的第二个主要定理。

\begin{theorem}{W}\label{thm:intro-W}
设 $R$ 是一个域。对任意有隙滤过 $A_\infty$ 代数 $(C,\mathfrak m_*)$，在
\[
  H(\overline C,\overline{\mathfrak m}_1)
  \otimes_R\Lambda_{0,\mathrm{nov}}(R)
\]
上都存在一个有隙滤过 $A_\infty$ 结构，它同伦等价于 $(C,\mathfrak m_*)$。
\end{theorem}

定理 W 是同调代数中一个经典结果的滤过 $A_\infty$ 代数版本；该结果可追溯至 \textbf{[Kad82]}。我们的证明采用对树求和的费曼图技术；它类似于 \textbf{[KoSo01]} 给出的方法，也类似于本书 2000 年版 \textbf{[FOOO00]} 第 A.6 节的论证。

第 6 章的目的是构造定理 D 中出现的谱序列，并说明它在辛几何与拓扑学中的各种应用。万有诺维科夫环 $\Lambda_{0,\mathrm{nov}}$ 带有滤过
\[
  F^\lambda\Lambda_{0,\mathrm{nov}}
  =T^\lambda\Lambda_{0,\mathrm{nov}}
\]
（见 (Conv.6)）。不过，这个滤过的参数 $\lambda$ 取实数，而不是整数。事实表明，从这样的滤过复形构造谱序列并证明其收敛性，是一个非平凡的问题。第 6 章解决这一问题并证明定理 E。为此，我们还使用算子 $\mathfrak p$，以及为给定拉格朗日子流形 $L$ 配置的滤过 $A_\infty$ 代数中幺元的存在性。

第 7 章讨论横截性问题。这是构造本书所用各类伪全纯圆盘模空间之（虚）基本链的两个主要要素之一。（另一个是定向问题，见第 8 章。）为此，我们系统地使用 \textbf{[FuOn99II]} 所引入的仓西结构\liftIndex{I037}及其多值截面。

对于所研究对象只有有限自同构群的一般模问题，仓西结构以及用多值截面作扰动，提供了研究横截性问题的最佳框架。尽管这套成熟工具自 1996 年 1 月首次出现以来已逾十年，相关领域的许多研究者似乎仍未充分注意到它。因此，我们在附录第 A1 节相当详细地叙述仓西结构，以及借助仓西结构构造（虚）基本链的方法。这里也想较为详细地介绍第 7 章的内容。

第 7.1 节在本书使用的各类伪全纯圆盘模空间上构造仓西结构。在定义 $A_\infty$ 代数和弗洛尔上同调的情形中，我们必须使用（虚）基本链，而不是（虚）基本闭链。这意味着，我们构造的各种代数结构之结构常数依赖于各种扰动的选择，因此只有它们的同伦型才可能成为不依赖这些扰动的不变量。由于必须考虑与拉格朗日子流形 $L$ 或拉格朗日子流形对 $(L_0,L_1)$ 相伴的一整族模空间，如果希望最终结果成为 $L$ 或 $(L_0,L_1)$ 本身的不变量，就必须为各个模空间选择满足某些相容条件的扰动。（这一点与通常的格罗莫夫--威腾不变量形成鲜明对比：后者在同调层面定义，因而可以放心地分别为每个模空间选择扰动。）我们在附录第 A1 节定义仓西结构的纤维积，并在第 7.1.1 小节用它来表述紧化模空间不同分层上所定义的仓西结构之间的相容性。第 7.2 节则借助这种相容性，仔细陈述多值截面（扰动）的相容条件。我们强调，如果不精确陈述扰动的相容条件，任何声称“利用紧化模空间定义的各种代数结构是良定义的”之陈述，都将失去意义和依据。

第 7.1.2--7.1.3 小节讨论我们所关心的圆盘模空间之仓西结构构造的解析部分。其中许多论证综合了文献中已经出现的各种技术。例如，博特--莫尔斯情形下的粘合技术见 \textbf{[Fuk96II]}，带模参数的伪全纯圆盘之研究见 \textbf{[FuOh97]}，仓西结构的粘合构造则见 \textbf{[FuOn99II]}。

用于定义拉格朗日子流形所配滤过 $A_\infty$ 结构的圆盘模空间之仓西结构构造，在第 7.1.4 小节完成。

第 7.1.5 小节研究一个模空间，我们用它定义拉格朗日子流形所配滤过 $A_\infty$ 代数之间的滤过 $A_\infty$ 同态。为了构造满足上述相容条件的仓西结构，需要研究圆盘 $D^n$ 的某种胞腔分解；反过来，这一胞腔分解又用来定义 $A_\infty$ 映射的概念（定理 7.1.51）。

注意，斯塔谢夫胞腔，更准确地说，斯塔谢夫对圆盘 $D^n$ 的胞腔分解，被用于定义 $A_\infty$ 空间的概念 \textbf{[Sta63]}；它既同带标记点的亏格零有边黎曼曲面模空间有关，也同度量带状树模空间有关（见 \textbf{[FuOh97]}）。斯塔谢夫告诉我们，他已经知道那种导向 $A_\infty$ 映射定义的圆盘胞腔分解（另见 \textbf{[MSS02, BoVo73]}）。不过，我们对这一胞腔分解的几何描述似乎是新的。

我们强调，本书关于横截性的叙述，其新意并不在解析部分——原则上，这一部分自 1996 年起便已得到理解——而在第 7.2 节给出的拓扑部分。在第 7.2 节，为了构造滤过 $A_\infty$ 结构等对象，我们在紧化模空间的每一分层上，对给定伪全纯方程的仓西映射逐层作适当的归纳扰动，并选择奇异余链复形的一个可数生成子复形。

\begin{remark}{1.32}\label{rem:1-32}
本书使用奇异链复形来构造拉格朗日子流形所配的 $A_\infty$ 结构。这是我们仔细斟酌后的选择。事实上，此外还有另外两种可能的选择；下面简述这两种选择，并说明本书为何采用现在的方案。

一种选择是使用德拉姆复形\liftIndex{I036}。正如 \textbf{[Fuk05II]} 所解释、\textbf{[Fuk07I, FOOO08II]} 所具体实行的那样，这一选择确实可行（这些工作预设了本书的内容）。它有自身的一些优点，例如在构造 $A_\infty$ 结构时更容易保持循环对称性等较多的对称性。不过，使用德拉姆复形便无法取 $\mathbb Z_2$ 系数，而 $\mathbb Z_2$ 系数对辛几何中的若干应用至关重要，尤其是对阿诺尔德--纪梵塔尔猜想的研究。我们将在 \textbf{[FOOO09I]} 中讨论本书理论对 $\mathbb Z$ 或 $\mathbb Z_2$ 系数的部分推广，并用它证明某些情形下的阿诺尔德--纪梵塔尔猜想；那里会采用本书的代数框架以及其他框架。因此，除第 7.5 节和定理 X 外，本书没有采用德拉姆复形。

另一种选择是使用莫尔斯复形。事实上，第一署名作者于 1993 年开创的正是这种方法（见 \textbf{[Fuk93]} 以及 \textbf{[Fuk97III]} 的 (3.3) 和图 8）。第二署名作者也采用这种方法构造了一个类似于定理 D 中谱序列的谱序列，并在 \textbf{[Oh96I, 96II]} 中指出，这种方法可以扩展为 $A_\infty$ 范畴\liftIndex{I006}的构造。本书 2000 年版第 287--291 页也用过类似论证来证明弱有限性（定义 6.3.27），所利用的是紧流形的莫尔斯复形总是有限生成这一显然事实。在准备当前修订版的过程中，我们为可数生成复形的弱有限性找到了一个（纯代数且更简单的）证明，因而选择收入这个较简单的证明，并从当前版本中删去旧证明。

我们相信，采用莫尔斯复形的方法在一定程度上可行。不过，这种方法会遇到一些麻烦。回忆与 $L$ 相伴的（滤过）$A_\infty$ 代数上的乘积结构，是上同调 $H^*(L)$ 上杯积的一个形变。第一署名作者 \textbf{[Fuk93,97III]} 以及 Martin Betz 与 Ralph L. Cohen \textbf{[BeCo94]} 的研究表明，为了在链层面研究莫尔斯同调的乘积，需要在一个流形上同时使用若干不同的莫尔斯函数。莫尔斯同调中的乘积，是通过考察从“Y-图”到给定流形的映射之模空间来定义的。为了定义结构常数，我们需要这个模空间具有横截性。注意，如果在 Y-图的三条边上采用同一个莫尔斯函数，这个模空间绝不会横截。这一横截性问题是代数拓扑中的一个基本问题；除非同时使用若干莫尔斯函数，否则似乎极难解决。参见 \textbf{[FOOO08III]}。

上述讨论的一个后果是，我们只能构造某种滤过 $A_\infty$ 范畴，其中的态射（即弗洛尔上同调）只能在一般位置下复合。这就导向拓扑 $A_\infty$ 范畴的概念。结果，要定义无非就是恒等态射的幺元便十分困难。第一署名作者在 \textbf{[Fuk97II]} 中于拓扑 $A_\infty$ 范畴的语境下讨论了幺元问题，但 \textbf{[Fuk97II]} 的叙述很不令人满意。只有在我们成功地按照本书采用的方式（或用德拉姆上同调）为每个 $L$ 构造（滤过）$A_\infty$ 代数之后，才得以正确理解（滤过）$A_\infty$ 范畴的幺元（见 \textbf{[Fuk02II]}）。

幺元在辛几何的应用中起着非常重要的作用，因为在某些有利情形下，可以用它证明弗洛尔上同调的非平凡性（例如见定理 E）。幺元在镜像对称的应用中同样重要；例如，$A_\infty$ 范畴的米田引理之证明以幺元的存在性为基础（见 \textbf{[Fuk02II]}）。因此，在这里正确地定义幺元，是拉格朗日子流形的 $A_\infty$ 范畴之一个重要组成部分。由于在使用莫尔斯函数的方法中似乎不易正确地定义幺元，本书也不使用莫尔斯同调。

即使完全不考虑（伪）全纯圆盘的贡献，横截性问题也已经存在。在经典代数拓扑中，例如，这个问题表现为：如何在奇异链复形上定义一个 $A_\infty$ 结构，使它在链层面实现交积。奇异链复形的主要困难在于，一条链绝不会与自身横截，因而同一条链与自身的配对并非良定义，除非这条链是底流形的基本链。代数拓扑学家熟知这一困难，也熟知在链层面实现交积的问题何其困难（除非使用德拉姆复形）。我们通过仔细扰动 $\mathbb Z$ 上通常的链交配对，解决了这一横截性问题；但所得乘积的结合律随后必然失效。我们系统地引入高阶乘积运算 $\overline{\mathfrak m}_k$ 来补偿结合律的失效，从而在一个可数生成的（余）链复形上得到 $A_\infty$ 代数。完成这一步之后，再以某种相容方式继续纳入伪全纯圆盘的贡献，并定义与给定拉格朗日子流形 $L$ 相伴的滤过 $A_\infty$ 代数。
\end{remark}

第 7.3 节构造同伦幺元。

第 7.5 节证明下面的定理。它很接近注 1.32 所提到的各种方法之间的等价性陈述；这些方法都用来构造定理 A 中出现的滤过 $A_\infty$ 代数。

回忆德拉姆复形是一个微分分次代数，因而可以视作一个 $A_\infty$ 代数。（这一说法的精确含义见注 3.2.4 (2)。）

\begin{theorem}{X}\label{thm:intro-X}
设 $(H^*(L;\Lambda_{0,\mathrm{nov}}),\mathfrak m_*)$ 是定理 A 给出的滤过 $A_\infty$ 代数，而 $(H^*(L;\mathbb Q),\overline{\mathfrak m}_*)$ 是将系数环约化到 $\mathbb Q$ 后所得的代数。则 $A_\infty$ 代数
\[
  (H^*(L;\mathbb Q),\overline{\mathfrak m}_*)
  \otimes_{\mathbb Q}\mathbb R
\]
作为 $A_\infty$ 代数同伦等价于 $L$ 的德拉姆复形。
\end{theorem}

这一定理蕴含：如果取 $\mathbb R$ 系数，那么定理 A 中的滤过 $A_\infty$ 代数是 $L$ 的实同伦型的一个形变。

直观上或许可以说，奇异链、德拉姆与莫尔斯这三种方法最终都会归结为同一套理论。然而，要严格证明这种等价并不容易。定理 X 给出了德拉姆上同调版本与奇异上同调版本之间等价性证明的一部分。很可能进一步精细化这些论证，再使用 \textbf{[Fuk07I]} 中构造的德拉姆版本，就能完整证明这一等价性。

至于经典拓扑中莫尔斯同调（或莫尔斯同伦）版本与德拉姆上同调版本之间的等价性，一个相关陈述是：\textbf{[FuOh97]} 所构造的给定紧流形之莫尔斯 $A_\infty$ 范畴，同伦等价于该流形的有理同伦型。正如我们在 \textbf{[FuOh97]} 中提到的，我们毫不怀疑这一陈述成立。遗憾的是，文献中尚未详细证明它。孔采维奇和 Yan Soibelman 在 \textbf{[KoSo01]} 第 239 页定理 2 附近的讨论，似乎与这一陈述有关；但 \textbf{[KoSo01]} 缺少该陈述的证明细节。（不过，他们在那里提出的证明思路颇具说服力，而且很可能是正确的。）

\begin{remark}{1.33}\label{rem:1-33}
我们猜想
\[
  \bigl(H^*(L;\mathbb Q),
    \{\overline{\mathfrak m}_k\}_{k=0}^{\infty}\bigr)
\]
同伦等价于 $\mathbb Q$-德拉姆复形（见 \textbf{[Sul78]}）。这一结论并不能由定理 X 推出。事实上，存在一对 $\mathbb Q$ 上的微分分次代数：它们在 $\mathbb R$ 上彼此同伦等价，却在 $\mathbb Q$ 上彼此不同伦等价（见 \textbf{[BrSz89]}）。
\end{remark}

第 7.4 节给出第 3.8 节所概述构造的细节。也就是说，我们利用同时带内部标记点和边界标记点的伪全纯圆盘模空间，形变与拉格朗日子流形相伴的滤过 $A_\infty$ 代数。

本节讨论两个主要问题。其一是只有内部标记点而没有边界标记点的伪全纯圆盘模空间之紧化。我们发现的一个实质性新现象是：标准的稳定映射紧化在这种情形下\emph{并不给出}正确的紧化。由于这一现象，刻画算子 $\mathfrak p$ 之主要性质的公式 (3.8.10.3) 中会出现含格罗莫夫--威腾不变量的项。

第 7.4 节处理的另一个问题，是为
$(H(L;\Lambda_{0,\mathrm{nov}}),\mathfrak m_*)$ 的 bulk 形变之描述提供一个代数框架。形变理论的一般方针（其源头可追溯至 \mbox{Murray Gerstenhaber}）断言，一个代数系统的形变由适当的李代数控制。对于滤过 $A_\infty$ 代数 $C$，有 Hochschild 链复形
\[
  CH(C,C)=\prod_k\operatorname{Hom}(B_kC[1],C[1]),
\]
它是一个微分分次李代数。我们把这个复形视作一个 $L_\infty$ 代数，并记作 $\operatorname{Der}B(C[1])$（见定义 7.4.26；$L_\infty$ 代数的定义见第 A3 节）。它的保幺版本记作 $\operatorname{Der}B^{\mathrm{unit}}(C[1])$。另一方面，把环境辛流形的上同调群 $H(M;\Lambda_{0,\mathrm{nov}})[1]$（次数已经移位）视作一个所有运算均平凡的 $L_\infty$ 代数。算子 $\mathfrak q$\liftIndex{I038} 的主要性质可以用 $L_\infty$ 代数的语言陈述如下。

\begin{theorem}{Y}\label{thm:intro-Y}
存在一个严格且滤过的 $L_\infty$ 同态
\[
  \mathfrak q^\circ\colon H(M;\Lambda_{0,\mathrm{nov}})[1]
  \longrightarrow
  \operatorname{Der}B^{\mathrm{unit}}
  \bigl(H(L;\Lambda_{0,\mathrm{nov}})[1]\bigr),
\]
使得把系数环约化到 $\mathbb Q$ 之后，映射 $\mathfrak q^\circ$ 约化为由包含映射在上同调群 $H(L;\mathbb Q)$ 上诱导的线性映射
\[
  H(M;\mathbb Q)
  \longrightarrow
  \operatorname{Hom}(\mathbb Q,H(L;\mathbb Q))
  \cong H(L;\mathbb Q).
\]

若 $\psi\colon(M,L)\to(M',L')$ 是辛微分同胚，则有下图；它在一个 $L_\infty$ 同伦的意义下交换：
\[
\begin{array}{ccc}
H(M;\Lambda_{0,\mathrm{nov}})[1]
  & \xrightarrow{\ \mathfrak q^\circ\ }
  & \operatorname{Der}B^{\mathrm{unit}}
    \bigl(H(L;\Lambda_{0,\mathrm{nov}})[1]\bigr) \\
\underset{\psi_*}{\Big\downarrow}
  && \underset{\psi_*}{\Big\downarrow} \\
H(M';\Lambda_{0,\mathrm{nov}})[1]
  & \xrightarrow{\ \mathfrak q^\circ\ }
  & \operatorname{Der}B^{\mathrm{unit}}
    \bigl(H(L';\Lambda_{0,\mathrm{nov}})[1]\bigr).
\end{array}
\]
\end{theorem}

定理 Y 是定理 7.4.118 和 7.4.120 的合并；这两个定理分别在第 7.4.5、7.4.6 小节证明。

利用 $L_\infty$ 模的概念，可以为映射 $\mathfrak r$ 陈述一个类似结果。这个算子出现在用环境辛流形的上同调类对拉格朗日子流形对之弗洛尔上同调所作的形变中；参见第 7.4.7--7.4.9 小节的定理 7.4.154、7.4.156、7.4.162。算子 $\mathfrak p$ 则可借助 $H(L;\Lambda_{0,\mathrm{nov}})$ 的循环 bar 复形来解释；这个复形带有一个以 $\operatorname{Der}B^{\mathrm{unit}}(H(L;\Lambda_{0,\mathrm{nov}})[1])$ 为作用代数的 $L_\infty$ 模结构。参见第 7.4.11 小节的定理 7.4.192 和 7.4.195。

第 7.4 节的结果与拉格朗日弗洛尔理论的环路空间表述 \textbf{[Fuk05II]} 密切相关，也与把我们的理论推广到高亏格情形密切相关。

第 8 章专门讨论我们构造中的另一个技术核心，即定向问题。第 8.1 节研究不带标记点的伪全纯圆盘模空间。这个问题自然地联系到椭圆算子族的指标理论。在全纯圆盘的情形（更一般地，在有边黎曼曲面的情形），我们需要带边界条件的椭圆算子族之指标理论。我们说明第二施蒂费尔--惠特尼类如何进入这一指标理论，并证明：若拉格朗日子流形相对 spin，则不带标记点的伪全纯圆盘模空间可定向。更准确地说，我们证明下面的定理 Z。以 $\mathcal M(L;\beta)$ 表示同调类 $\beta\in H_2(M,L)$ 中伪全纯圆盘
\[
  w\colon(D^2,\partial D^2)\longrightarrow(M,L)
\]
之模空间的稳定映射紧化。

\begin{theorem}{Z}\label{thm:intro-Z}
若 $L\subset(M,\omega)$ 是相对 spin 的拉格朗日子流形，则模空间 $\mathcal M(L;\beta)$ 可定向。此外，$L$ 上相对 spin 结构的选择，会对所有 $\beta\in H_2(M,L)$ 的 $\mathcal M(L;\beta)$ 典范地确定一族相容定向。
\end{theorem}

相对 spin 结构的定义见第 8 章定义 8.1.2。还要指出，Vin de Silva 在 \textbf{[Sil97]} 中独立证明了定理 Z。

为了说明带边界条件的柯西--黎曼算子之指标丛一般并不可定向，我们给出一个具体例子：圆盘上带全实边界条件的柯西--黎曼算子之指标丛不可定向。

第 8 章其余各节致力于定义各种伪全纯圆盘模空间关于纤维积的相容定向。这个问题事实证明是构造 $A_\infty$ 代数时的一个要点。

给出相容定向的问题，同为 $A_\infty$ 公式寻找合适符号约定的问题紧密相连，也与超对称密切相关。按照 Ezra Getzler 与 John D. S. Jones \textbf{[GeJo90]} 的说法，$A_\infty$ 公式基本上有两类不同的符号约定。下面简要回顾它们的差别。考虑 $L$ 的（奇异）链复形 $C(L)$。设 $x_i\in C$，并令其次数为
\[
  d_i=\operatorname{codim}x_i.
\]

在第一种约定中，我们把次数移位一\liftIndex{I028}，并在移位后的复形中把它视为次数 $d_i+1$ 的元素。依照这种约定，在乘积中交换两个元素时产生的符号变化为
\[
  (-1)^{(d_1+1)(d_2+1)}.
\]

在第二种约定中，每个元素有两种不同的次数：一种是分次复形中的通常次数，另一种是超次数。在这种约定下，元素 $x_i\in C(L)$ 的通常次数是分次给出的 $d_i$，超次数则是 $1$。在乘积中交换两个元素时产生的符号变化为
\[
  (-1)^{d_1d_2+1}.
\]

在这个问题上，Getzler--Jones 以 $\widetilde{\mathfrak m}$ 表示采用第一种符号约定的算子，以 $\mathfrak m$ 表示采用第二种符号约定的算子。二者由 \textbf{[GeJo90]} 第 259--260 页的映射 $b$ 联系起来。本书的约定似乎与第一种符号约定一致。（注意，\textbf{[Fuk02III]} 采用的是第二种符号约定。）

从几何观点看，这两种符号约定对应于两种不同的方式，用来给边界上带标记点的伪全纯圆盘模空间赋予定向。

因此，第 8.2--8.10 节详细描述各种模空间的定向，并用第一种符号约定验证带符号的 $A_\infty$ 公式。我们如此详尽地说明这个定向问题，是因为它既复杂又容易引起混淆。若无彻底分析，便无法看清贯穿全书的正确带符号公式。
